\chapter{The Lean Target}
\label{ch:lean}

This chapter is a tour of the generated \Lean{}~4 certificate.  The
files of interest are:

\begin{itemize}[leftmargin=1.7em]
  \item \texttt{deppy/lean/sidecar\_certificate.py} ---
    \texttt{render\_sidecar\_section}, the orchestrator;
  \item \texttt{deppy/lean/sidecar\_annotation\_preamble.py} ---
    the concrete-Props preamble;
  \item \texttt{deppy/lean/sidecar\_mechanization.py} ---
    Python-axiom $\to$ Lean-Prop translation;
  \item \texttt{deppy/lean/sidecar\_type\_inference.py} ---
    \texttt{inspect.signature} $\to$ Lean signatures;
  \item \texttt{deppy/lean/sidecar\_body\_link.py} ---
    body translation per linked method;
  \item \texttt{deppy/lean/body\_translation.py} ---
    AST $\to$ Lean expression;
  \item \texttt{deppy/lean/sidecar\_derivation.py} ---
    legacy \texttt{by:}-clause derivations.
\end{itemize}

We give the certificate's structure top-to-bottom and explain how the
trust chain bottoms.

\section{Trust chain in the Lean output}
\label{sec:trust-chain}

The certificate is laid out so the trust chain is visible:

\begin{equation*}
  \boxed{\text{Foundations}} \;\longrightarrow\; \boxed{\text{Sidecar axioms}} \;\longrightarrow\; \boxed{\text{Lemmas}} \;\longrightarrow\; \boxed{\text{Theorems / Verify blocks}}.
\end{equation*}

Each box represents a contiguous section in the \Lean{} file:
\begin{description}
  \item[Foundations] are the bottom: arithmetic identities admitted
  on the user's authority (or proved by \texttt{omega} when Z3
  succeeds).  Foundations have no upstream dependencies in the
  certificate.
  \item[Sidecar axioms] depend on foundations: each axiom is a
  claim about library behaviour, expressed as a concrete \Lean{}
  Prop using the mechanization vocabulary.
  \item[Lemmas] depend on foundations and axioms.  When the user
  supplied \texttt{lean\_proof:}, lemmas are real \texttt{theorem}s.
  \item[Theorems] are the top: derived from any combination of the
  above.  Verify blocks emit one \texttt{theorem} each, with the
  proof body being the PSDL-compiled tactic block (or a derived
  proof from a \texttt{by:}-clause for legacy proofs).
\end{description}

\Lean's elaborator processes the file top to bottom, so dependent
declarations resolve.  The compiler counts how many declarations
end up at each level (foundations, axioms, lemmas, theorems), how
many of each carry concrete bodies vs. opaque placeholders, and
emits these in the certificate's header comment.

\section{Header}

Every certificate begins with:

\begin{lstlisting}[language=lean4]
-- Auto-generated by ``deppy.lean.sidecar_certificate``.
-- Module: Deppy_<sidecar_basename>
-- Foundations: <N>; specs: <M>; axioms: <K>; proofs: <P>.
-- Library: sympy 1.14.0 -- sympy
--
-- Trust chain: Foundations (mathematical facts) -> Sidecar axioms
-- (claims about library behaviour) -> Theorems (real proof
-- terms -- Eq.trans / And.intro / rfl / Foundation.symm.trans
-- -- type-checked by Lean's kernel).
\end{lstlisting}

\section{Annotation-driven concrete preamble}
\label{sec:concrete-preamble}

The first substantive section is the annotation-driven concrete
preamble (file: \texttt{deppy/lean/sidecar\_annotation\_preamble.py}).
Its purpose is to replace the \emph{opaque}
\texttt{Sidecar\_X\_prop : Prop} / \texttt{Foundation\_X\_prop : Prop}
declarations of older versions with \emph{concrete} \Lean{}
propositions assembled from the sidecar's \texttt{binders:},
\texttt{function:}, \texttt{property:}, and \texttt{code\_types:}
fields.

The motivation, in \texttt{AUDIT.md}'s words: ``Once Props are
concrete, the trust chain shifts from `\Deppy{} admitted opaque
X' to `\Lean's kernel checked the user's PSDL tactic against a
concrete claim'.''

\subsection{Concrete types}

For each unique binder type across all verify blocks, the preamble
emits an \texttt{abbrev}:

\begin{lstlisting}[language=lean4]
-- Concrete types from verify-block binders (definitional aliases to Int)
abbrev Circle   : Type := Int
abbrev Point    : Type := Int
abbrev Polygon  : Type := Int
abbrev Segment  : Type := Int
abbrev Triangle : Type := Int
\end{lstlisting}

The \texttt{abbrev T : Type := Int} is a deliberate flat-Int
simplification.  We acknowledge it openly: \emph{the generated
certificate identifies every Python class with the carrier
\texttt{Int}}.  This loses the ability to distinguish classes in
\Lean's type system but lets the certificate compile without
class-machinery dependency.  In a future iteration where the
preamble engages the \texttt{deppy.lean.sidecar\_type\_inference}
result, types could be richer (e.g. structured records or
quotient types over the underlying Python instance representation).

\subsection{Concrete code-type signatures}

For each entry in the sidecar's top-level \texttt{code\_types:}
block:

\begin{lstlisting}[language=lean4]
-- User-declared code-type signatures (from .deppy code_types:)
axiom sqrt           : Int → Int
axiom sum_zip_sub_sq : Point → Point → Int
\end{lstlisting}

These are the helpers PSDL references but the library does not
necessarily expose at the top level.  The \texttt{axiom} kind means
they are uninterpreted; the \PSDL-compiled tactics use them as
applied symbols, and the kernel does not unfold them.

\subsection{Concrete foundation Props}

Each foundation gets a concrete \texttt{def} or an \texttt{opaque}
fallback, depending on whether \texttt{discharge\_foundations} was
able to verify it via Z3 or whether the user supplied an explicit
\texttt{lean\_proof}.

For arithmetic foundations Z3 verifies:
\begin{lstlisting}[language=lean4]
-- Foundation: Real_add_comm
-- Statement: a + b == b + a
def Foundation_Real_add_comm_prop : Prop :=
  ∀ (a b : Int), a + b = b + a
theorem Foundation_Real_add_comm : Foundation_Real_add_comm_prop := by
  unfold Foundation_Real_add_comm_prop; intros; omega
\end{lstlisting}

For foundations involving \texttt{sqrt}:
\begin{lstlisting}[language=lean4]
-- Foundation: Real_sqrt_nonneg
-- Statement: sqrt(x) >= 0 when x >= 0
opaque Foundation_Real_sqrt_nonneg_prop : Prop
axiom Foundation_Real_sqrt_nonneg : Foundation_Real_sqrt_nonneg_prop
\end{lstlisting}

The certificate's foundation section is sorted alphabetically and
groups Z3-verified foundations (real \texttt{theorem}s) above
admitted ones (\texttt{axiom}s).  The classification feeds the
report's \texttt{foundations\_z3\_verified} counter.

\section{Mechanization preamble}
\label{sec:mechanization-preamble}

After the annotation-driven preamble, the certificate emits the
\emph{mechanization vocabulary}: the union of types, methods,
attributes, and helpers extracted from grounded axiom AST walks.

\begin{lstlisting}[language=lean4]
-- Shared attribute accessors (derived from var.attr)
axiom dot_x        : Int → Int
axiom dot_y        : Int → Int
axiom dot_centroid : Int → Int
axiom dot_radius   : Int → Int
-- ... and so on

-- Class-prefixed methods and constructors
axiom Triangle_mk : Int → Int → Int → Int
axiom Segment_mk  : Int → Int → Int
axiom Circle_mk   : Int → Int → Int
def Point_distance : Int → Int → Int := fun _0 _1 => 0
def Polygon_area  : Int := 0

-- Free functions
axiom abs           : Int → Int
axiom area          : Int → Int
axiom is_collinear  : Int → Int → Int → Prop
axiom thales        : Int → Int → Int → Int → Prop
\end{lstlisting}

The methods that have a body translation (see
\texttt{body\_translation}) are emitted as \texttt{def := sorry}
\emph{placeholders} so that \texttt{unfold} in proof bodies can
fire without elaborating the full body inline.  The actual
translated body is emitted in commented form below the
\texttt{Body links} section.  This dual encoding (placeholder
\texttt{def} + commented body) preserves elaboration-friendliness
while documenting what the body actually translates to.

\subsection{Body links section}

For each method whose body the translator could read, the
certificate emits:
\begin{lstlisting}[language=lean4]
-- ── Point_distance ── (sympy.geometry.point.Point.distance)
-- Source SHA-256: df5da0ce57c1b906...
-- deppy body-translator: EXACT (full translation, no sorry)
-- Lean (commented to avoid collision with opaque preamble decl):
-- def Point_distance (self : Int) (other : Int) : α :=
--   if (isinstance other Point) then
--     let (s, p) := (Point._normalize_dimension self (Point other));
--     (sqrt (Add (List.zipWith (fun a b => ((a - b) ^ 2)) s p)))
--   else
--     let distance := (getattr other "distance" ());
--     if (distance = ()) then panic! "exception" else (distance self)
\end{lstlisting}

The hash is the SHA-256 of the inspected source bytes, providing a
commitment to which version the certificate verified against.  When
the source changes, the hash changes; the
\texttt{expects\_sha256:} field on \texttt{verify} blocks can pin a
specific hash and the \texttt{certify\_or\_die} gate fires
\texttt{hash\_drift\_count} on mismatch.

\section{Sidecar axioms section}
\label{sec:axioms-section}

For each \texttt{AxiomDecl}, the certificate emits one
\texttt{Sidecar\_X\_prop} \texttt{def} (when grounded and
mechanizable) or one \texttt{opaque} (otherwise), followed by the
\texttt{axiom Sidecar\_X : Sidecar\_X\_prop} declaration.

\begin{lstlisting}[language=lean4]
-- Axiom: dist_pythagoras
-- Grounding: GROUNDED
-- Target: sympy.geometry.point.Point; Module: sympy
-- Statement (Python): Point.distance(p, q)**2 == (p.x - q.x)**2 + (p.y - q.y)**2
-- Concrete Lean Prop (mechanized from .deppy):
def Sidecar_dist_pythagoras_prop : Prop :=
  ∀ (p q : Int),
    (Point_distance p q) ^ 2 =
      ((dot_x p) - (dot_x q)) ^ 2 + ((dot_y p) - (dot_y q)) ^ 2
axiom Sidecar_dist_pythagoras : Sidecar_dist_pythagoras_prop
\end{lstlisting}

For non-grounded axioms (target missing, method missing, or
natural-language documentation):
\begin{lstlisting}[language=lean4]
-- Axiom: euler_line_collinear
-- Grounding: METHOD_MISSING
-- ... Cannot mechanize: <reason>
opaque Sidecar_euler_line_collinear_prop : Prop
axiom Sidecar_euler_line_collinear : Sidecar_euler_line_collinear_prop
\end{lstlisting}

The certificate's \texttt{nl\_axioms},
\texttt{target\_missing\_axioms}, \texttt{method\_missing\_axioms},
and \texttt{mechanized\_axioms} counters are populated by this pass.

\section{Verify blocks: theorem emission}
\label{sec:verify-emission}

For each successful verify block (kernel and PSDL-compile both
returned success), the certificate emits a tactic-mode \Lean{}
theorem:

\begin{lstlisting}[language=lean4]
-- Verify block: dist_nonneg_psdl
-- Function: sympy.geometry.point.Point.distance
-- Property: Point.distance(p, q) >= 0
-- Foundation: Real_sqrt_nonneg
-- Source SHA-256: df5da0ce57c1b906...
-- Trust level: STRUCTURAL_CHAIN
-- Axioms used: ['Real_sqrt_nonneg']
-- Binders: p: Point, q: Point

def Verified_dist_nonneg_psdl_property : Prop :=
  ∀ (p q : Point), Point_distance p q ≥ 0

theorem Verified_dist_nonneg_psdl :
    Verified_dist_nonneg_psdl_property := by
  intros p q
  -- PSDL compiled output:
  by_cases h : isinstance other Point
  · -- value-fibre: isinstance(other, Point)
    unfold Point_distance  -- expected: sqrt(sum_zip_sub_sq(self, other))
    exact Foundation_Real_sqrt_nonneg
    have h_493 : sum_zip_sub_sq self other ≥ 0 := by omega
      -- z3-discharged
  · -- complement-fibre: ¬ isinstance(other, Point)
    -- vacuous: NotImplementedError
    exact absurd ‹_› (by simp)
\end{lstlisting}

This is the pay-off.  Reading top-down:
\begin{enumerate}[leftmargin=1.7em]
  \item \texttt{Verified\_dist\_nonneg\_psdl\_property : Prop} is a
    \emph{concrete} Lean proposition whose body is the universal
    quantification of the property over the binders.  This is what
    \texttt{build\_annotation\_preamble} produces.
  \item The theorem's goal type is exactly that \texttt{Prop}.
  \item The proof body is the PSDL-compiled tactic block, with
    \texttt{intros p q} prepended (because the property quantifies
    over $p, q$).
  \item \Lean's elaborator must close the body against the goal,
    using the foundations, axioms, and code-types declared above.
\end{enumerate}

When \Lean{} accepts, this is an honest verification.  When it
rejects, the \texttt{lean\_runner.stderr} is captured and printed,
and \texttt{certify\_or\_die} returns failure.

\subsection{Failure modes for verify-block elaboration}

Three classes of failure are common in practice:

\paragraph{Goal-type mismatch.}  The PSDL tactic emitted
\verb|exact Foundation_Real_sqrt_nonneg|, which has type
\texttt{Foundation\_Real\_sqrt\_nonneg\_prop} (a $\forall$).  But the
local goal after \texttt{by\_cases} and \texttt{unfold} is
\texttt{(sqrt (\dots)) ≥ 0}, not the foundation's full
$\forall$-statement.  Fix: PSDL should emit
\verb|exact Foundation_Real_sqrt_nonneg <args>|.  This is a known
gap; the current \texttt{apply()} sugar emits the bare
\texttt{exact} and relies on \Lean's elaborator to instantiate the
$\forall$, which works \emph{sometimes}.

\paragraph{Unfold has no body.}  The mechanization preamble's
\texttt{Point\_distance} is a \texttt{def := fun \_0 \_1 => 0}; the
\texttt{unfold} produces \texttt{0}, not the actual translated body.
The expected body lives in the commented-out section below.  When
\texttt{unfold}'s opening of \texttt{Point\_distance} is followed by
\texttt{exact Foundation\_Real\_sqrt\_nonneg}, \Lean{} cannot
unify $0 \ge 0$ with $\sqrt{x} \ge 0$ unless the rewriting tactic is
strong enough.

\paragraph{Class-method elaboration ambiguity.}  When the property
mentions \texttt{Point.distance(p, q)}, the renderer produces
\texttt{(Point\_distance p q)}.  When the binders' types
\texttt{Point := Int} (flat-Int simplification), this elaborates
trivially.  When a future version uses richer types, the
elaborator may face a class-resolution choice and need explicit
type ascription.  We do not currently handle this case explicitly.

These failure modes are the practical consequences of the flat-Int
simplification.  The next iteration of the pipeline (see
\texttt{AUDIT.md} sections B and F) will need richer typed
preamble emission.  For now, the certificate emits both the
PSDL-compiled tactic block \emph{and} the kernel's verdict, so a
user can check the kernel succeeded even when \Lean{} rejects.

\section{Legacy proofs section}

For each \texttt{ProofDecl} (the legacy form), the certificate runs
\texttt{derive\_all} and emits one of:

\begin{lstlisting}[language=lean4]
-- Proof: distance_pythagoras_proof
-- Target: sympy.geometry.point.Point; By: rewrite Real_sqrt_sq_nonneg then cite dist_pythagoras
-- Classification: REWRITE_THEN_CITE
theorem Proof_distance_pythagoras_proof : Sidecar_dist_pythagoras_prop := by
  rw [Foundation_Real_sqrt_sq_nonneg]
  exact Sidecar_dist_pythagoras
\end{lstlisting}

The \texttt{Classification} is one of the \texttt{ProofClassification}
enum values:
\texttt{REFL}, \texttt{FOUNDATION}, \texttt{AXIOM}, \texttt{CHAIN},
\texttt{REWRITE\_THEN\_CITE}, \texttt{COMPOSITION},
\texttt{TARGET\_MISSING}, \texttt{UNCLASSIFIED}.  The certificate's
\texttt{classification\_counts} dict counts these.

\paragraph{Cited proofs and ``true derivations''.}  The pipeline
draws a line between:
\begin{itemize}[leftmargin=1.7em]
  \item \emph{Cited proofs} --- the proof body is a renaming of an
    existing axiom (\texttt{by: axiom X} reduces to
    \texttt{exact Sidecar\_X}).  These contribute no new
    information; the certificate counts them in \texttt{cited\_count}.
  \item \emph{Derived proofs} --- the body is a real composition
    that produces something new (\texttt{by: chain F1, F2, F3} or
    \texttt{rewrite F then cite A}).  Counted in
    \texttt{derived\_count}.
\end{itemize}

A high \texttt{cited\_count} on a sidecar means the sidecar mostly
\emph{re-states} library claims rather than proving them; a high
\texttt{derived\_count} means the certificate is doing real
mathematical work.  We expose both honestly.

\section{Cubical and cohomology sections}

The certificate also emits:
\begin{itemize}[leftmargin=1.7em]
  \item A \texttt{Cubical analysis} section listing per-method
    cell counts, Kan candidates, and Euler characteristics.
  \item A \texttt{Cohomology cocycles} section with the C⁰/C¹/C²
    cocycle theorems (each typically of the form
    \texttt{theorem deppy\_c0\_<fn> : True := trivial} when the
    safety predicates are absent).
  \item A \texttt{Bicomplex} section with the bicomplex's edge,
    face, and Euler counts.
  \item A \texttt{Fibration descent} section with per-fibre
    verification outcomes.
\end{itemize}

These are diagnostic in the current implementation: the kernel
does not consult them when verifying a verify-block's \texttt{ProofTerm}.
\Cref{ch:cohomology} discusses what they buy and what they do
not.

\section{Counterexample-rejected theorems}

When \texttt{run\_counterexample\_search} finds a counterexample
for a foundation, every theorem citing that foundation is
\emph{replaced} with a rejection notice:

\begin{lstlisting}[language=lean4]
-- Theorem: distance_nonneg_proof
-- REJECTED: cited foundation Real_sqrt_nonneg has counterexample
--   x = -1 satisfies the negation
-- (would have emitted: theorem ... := exact Foundation_Real_sqrt_nonneg)
\end{lstlisting}

The \texttt{rejected\_theorem\_count} counter feeds into
\texttt{certify\_or\_die}.  This is the cleanest implementation of
``don't lie about a proof when its foundation is provably false''.

\section{Body translation in detail}
\label{sec:body-translation-detail}

We owe more detail on \texttt{body\_translation.translate\_body}, the
function that turns a Python AST into a Lean expression.  The
implementation is in \texttt{deppy/lean/body\_translation.py}.

\subsection{Two-pass design}

\paragraph{Pass 1: type inference.}
\texttt{infer\_locals\_types(fn\_node, param\_types=...)} returns a
\texttt{\{name: type\_text\}} dict.  Type evidence comes from:
\begin{enumerate}[leftmargin=1.7em]
  \item Annotated parameters (\texttt{def f(xs: list[int])}).
  \item Operations: \texttt{xs.append(\dots)} forces
    \texttt{xs : list}; \texttt{len(xs)} forces \texttt{xs : list /
    dict / str / tuple}; \texttt{xs[k] = v} forces \texttt{xs : list
    / dict}; arithmetic propagates numeric types.
  \item Assignment: \texttt{x = e} propagates \texttt{type(e)}.
  \item Return: the function's return type unifies with all
    \texttt{return} targets.
\end{enumerate}
The walker is best-effort: when the inference produces a mostly-empty
type-context, the translation falls back to \texttt{α} polymorphic
return types.

\paragraph{Pass 2: translation.}
A \texttt{\_BodyT} walker uses \texttt{locals\_types} to:
\begin{itemize}[leftmargin=1.7em]
  \item pick the right return type for the function's
    \texttt{def} signature;
  \item translate \texttt{len(xs)} to \texttt{xs.length} or
    \texttt{xs.size} depending on the inferred type;
  \item route list-indexing through \texttt{List.get?} /
    \texttt{Option.get!};
  \item select the right comparison: \texttt{=} for
    \texttt{Int}/\texttt{Bool}, \texttt{==} for non-decidable types.
\end{itemize}

\subsection{Statements}

\begin{itemize}[leftmargin=1.7em]
  \item \texttt{If}: \texttt{if cond then t else e}; if either branch
    diverges or returns, that branch's type is \texttt{α} polymorphic.
  \item \texttt{Return}: the body's value is the returned expression.
  \item \texttt{Assign} / \texttt{AnnAssign}: \texttt{let x := e}.
  \item \texttt{For}: \texttt{(range n).foldl (fun acc i => body) acc} or
    \texttt{xs.foldl} when iterating a list.
  \item \texttt{While}: not currently supported; emits \texttt{sorry}.
  \item \texttt{Try}: not supported; emits \texttt{sorry}.
  \item \texttt{Raise}: \texttt{panic! "exception"}.
\end{itemize}

\subsection{Expressions}

\begin{itemize}[leftmargin=1.7em]
  \item Arithmetic, comparison: direct.
  \item \texttt{abs(x)}: \texttt{x.natAbs} (cast to \texttt{Int}).
  \item Method calls \texttt{x.m(args)}: \texttt{(x.m args)}.
  \item Class method calls \texttt{C.m(args)}: \texttt{(C\_m args)}.
  \item Attribute access: \texttt{(dot\_attr x)}.
  \item Tuple: \texttt{Pair2\_mk a b} / \texttt{Pair3\_mk a b c}.
  \item List literal: \texttt{[a, b, c]} → \texttt{[a, b, c]} (Lean
    list).
  \item Comprehension: best-effort; \texttt{[f x | x in xs]} →
    \texttt{xs.map (fun x => f x)}; with filters → \texttt{filterMap}.
\end{itemize}

\subsection{Worked example: \texttt{Polygon.area}}

\Sympy's \texttt{Polygon.area} (truncated):

\begin{lstlisting}[language=psdl]
@property
def area(self):
    area = 0
    args = self.args
    for i in range(len(args)):
        x1, y1 = args[i - 1].args
        x2, y2 = args[i].args
        area += x1 * y2 - x2 * y1
    return simplify(area) / 2
\end{lstlisting}

The body translator produces (from the actual generated
\texttt{sympy\_geometry.lean}):

\begin{lstlisting}[language=lean4]
def Polygon_area (self : Int) : α :=
  let area := 0;
  let args := (dot_args self);
  let area := ((range (args.length)).foldl
    (fun area i =>
      let (x1, y1) := (dot_args (args.get! ((i - 1).toNat)));
      let (x2, y2) := (dot_args (args.get! (i.toNat)));
      (area + ((x1 * y2) - (x2 * y1)))) area);
  ((simplify area) / 2)
\end{lstlisting}

The translator caught the index arithmetic
(\texttt{i - 1} → \texttt{(i - 1).toNat}), the tuple unpack
(\texttt{x1, y1 = args[i - 1].args} via \texttt{dot\_args} +
destructuring), and the closing accumulator.  The return type stays
\texttt{α} because the divisor introduces a non-integer division.
This is honest --- the certificate flags the polymorphic return as
``\texttt{not exact}'' and the verify block targeting
\texttt{Polygon.area} sees the output and downgrades trust
accordingly.

\section{Type inference from \texttt{inspect.signature}}
\label{sec:type-inference}

\texttt{deppy/lean/sidecar\_type\_inference.py} bridges the live
Python signatures into the Lean preamble.  For each linked method:
\begin{enumerate}[leftmargin=1.7em]
  \item Call \texttt{inspect.signature(method)}.
  \item Extract each parameter's annotation; resolve to a Python
    type (or fall back to \texttt{Int}).
  \item Translate the type to a Lean type via
    \texttt{deppy.lean.type\_translation}.
  \item Build a Lean signature \texttt{(arg1 : T1) → \dots → R}.
\end{enumerate}

The translator handles:
\begin{itemize}[leftmargin=1.7em]
  \item \texttt{int}, \texttt{float}, \texttt{bool}, \texttt{str},
    \texttt{None}: direct.
  \item \texttt{list[T]} / \texttt{List[T]}: \texttt{List \rebr{T}}.
  \item \texttt{dict[K, V]}: \texttt{Std.HashMap \rebr{K} \rebr{V}}.
  \item \texttt{tuple[T1, T2, \dots]}: \texttt{T1 × T2 × \dots}.
  \item \texttt{Optional[T]}: \texttt{Option \rebr{T}}.
  \item \texttt{Union[T1, T2]}: \texttt{Sum \rebr{T1} \rebr{T2}}.
  \item Class types: a flat-\texttt{Int} alias (the simplification
    of \S\ref{sec:concrete-preamble}).
\end{itemize}

\section{Lean target: a complete picture}
\label{sec:lean-complete}

To close the chapter we trace through the structure of the
generated \texttt{sympy\_geometry.lean} (1308 lines) section by
section:

\begin{enumerate}[leftmargin=1.7em]
  \item Lines 1--9: header comment with module name and
    provenance.
  \item Lines 11--33: annotation-driven concrete preamble
    (5 \texttt{abbrev}s for class types, 2 \texttt{axiom}s for
    code-types).
  \item Lines 35--136: mechanization preamble
    (14 \texttt{axiom}s for shared accessors, 9 for class methods
    and constructors, 11 for free functions).
  \item Lines 137--216: foundations (14 declarations: 7 real
    \texttt{theorem}s discharged by Z3, 7 admitted as
    \texttt{axiom}s).
  \item Lines 218--525: sidecar axioms (38 declarations, all but a
    handful mechanized to concrete \texttt{def}s).
  \item Lines 526--850: legacy proofs (45 \texttt{theorem}s with
    bodies derived from \texttt{by:}-clauses).
  \item Lines 851--1100: verify-block theorems (11 blocks, varying
    level of detail).
  \item Lines 1101--1180: cubical, cohomology, bicomplex sections.
  \item Lines 1181--1308: counterexample-rejected theorems (none
    in this run), audit appendices, ProofTerm-zoo probes.
\end{enumerate}

The total is 1308 lines of \Lean{}~4 source for a sidecar with 14
foundations, 5 specs, 38 axioms, and 11 verify blocks.

\section{Summary}

The \Lean{}~4 target:
\begin{itemize}[leftmargin=1.7em]
  \item is a single self-contained \texttt{.lean} file with explicit
    trust chain;
  \item uses concrete types (currently flat-\texttt{Int} aliases),
    concrete code-type signatures, and concrete foundation /
    sidecar-axiom Props where the user's annotations permit;
  \item is structured top-to-bottom from foundations to theorems
    so \Lean's elaborator processes dependent declarations in
    order;
  \item emits one \texttt{theorem} per successful verify-block,
    with the proof body being PSDL-compiled tactic text;
  \item flags every \emph{cited} (renaming) vs.\ \emph{derived}
    (real) proof so the user can audit triviality;
  \item integrates body translation behind comments to avoid
    colliding with the opaque preamble \texttt{def}s;
  \item rejects any theorem whose foundation a counterexample
    refuted.
\end{itemize}

The next chapter dives deeper into the \Zthree{} backend: which
formulas it discharges, how counterexample search runs, and what
the trust promotion from \texttt{AXIOM\_TRUSTED} to
\texttt{Z3\_VERIFIED} actually means.
